\section{Status-Specific Snapshots}

RCOUNT treats status as part of the summary key. A contest summary is not just
``the total for P-001''; it is the total for a contest, reporting unit, optional
batch, and status. This lets an unofficial total and a canvassed total coexist
in the same package without overwriting one another.

\begin{center}
\small
\begin{tabularx}{\textwidth}{lXX}
\toprule
Status & Example source & Mechanical checks \\
\midrule
Unofficial & Election-night export & Contest arithmetic, jurisdiction rollup for unofficial summaries, source hash. \\
Canvassed & Canvass report or corrected export & Contest arithmetic, jurisdiction rollup for canvassed summaries, status-event evidence. \\
Recounted & Recount report & Same arithmetic on the recount snapshot plus recount event lineage. \\
Amended & Amended statement or court order & Arithmetic on amended snapshot plus public event reference. \\
Certified & Certification artifact & Consistency of records marked certified, not legal validity. \\
\bottomrule
\end{tabularx}
\end{center}

The jurisdiction-total check is status-specific. If a package contains both
unofficial and canvassed snapshots, the unofficial jurisdiction total must equal
the unofficial precinct summaries, and the canvassed jurisdiction total must
equal the canvassed precinct summaries. This prevents a verifier from comparing
records across statuses as if they belonged to one immutable table.

